% R567: fuente francesa intro-fr.tex:70--93, pp. 7--8.
Queda fuera de nuestra competencia ofrecer en esta introducción una reseña histórica, siquiera sumaria, de las nociones y los resultados expuestos.
El texto solo contendrá referencias que se consideren particularmente útiles para su comprensión, y únicamente indicaremos el origen de los resultados más importantes.
Al menos desde el punto de vista formal, los temas tratados en nuestra obra son bastante nuevos; esto explica la escasez de referencias a los padres de la geometría algebraica del siglo~\textsc{xix} y de comienzos del~\textsc{xx}, cuyos trabajos solo conocemos de oídas.
Conviene, no obstante, decir aquí unas palabras sobre las obras que han influido más directamente en los autores y contribuido al desarrollo del punto de vista de los esquemas.
Hay que citar en primer lugar el trabajo fundamental (FAC) de J.-P.~Serre, que ha servido de introducción a la geometría algebraica a más de un joven adepto (entre ellos uno de los autores de este tratado) desalentado por la aridez de los clásicos \emph{Foundations} de A.~Weil~\cite{I-18}.
Allí se demuestra por primera vez que la «topología de Zariski» de una variedad algebraica «abstracta» es perfectamente adecuada para aplicarle ciertas técnicas de la topología algebraica y dar lugar, en particular, a una teoría cohomológica.
Además, la definición de variedad algebraica que allí se da es la que se presta más naturalmente a la extensión de esta noción que desarrollamos aquí\footnote{Como nos ha señalado J.-P.~Serre, conviene observar que la idea de definir la estructura de variedad mediante un haz de anillos se debe a H.~Cartan, quien tomó esta idea como punto de partida de su teoría de los espacios analíticos.
Naturalmente, al igual que en geometría algebraica, sería importante dar cabida en «geometría analítica» a los elementos nilpotentes en los anillos locales de los espacios analíticos.
H.~Grauert~\cite{I-5} ha abordado recientemente esta extensión de la definición de H.~Cartan y J.-P.~Serre, y cabe esperar que pronto aparezca una exposición sistemática de geometría analítica en este marco general.
Por otra parte, es evidente que las nociones y las técnicas desarrolladas en este tratado conservan su sentido en geometría analítica, aunque haya que esperar dificultades técnicas mayores en esta última teoría.
Puede preverse que la geometría algebraica, por la sencillez de sus métodos, podrá servir como una especie de modelo formal para desarrollos futuros de la teoría de los espacios analíticos.}.
El propio Serre había observado, además, que la teoría cohomológica de las variedades algebraicas afines podía transcribirse sin dificultad sustituyendo las álgebras afines sobre un cuerpo por anillos conmutativos cualesquiera.
Los capítulos I y II de este tratado y los dos primeros párrafos del capítulo III pueden considerarse, por tanto, esencialmente como transposiciones sencillas, a este marco ampliado, de los resultados principales de (FAC) y de un artículo posterior del mismo autor~\cite{I-15}.
También hemos sacado gran provecho del \emph{Séminaire de Géométrie algébrique} de C.~Chevalley~\cite{I-1}; en particular, el uso sistemático de los «conjuntos construibles» que introdujo se ha revelado muy útil en la teoría de esquemas (véase el cap.~IV).
Le hemos tomado también el estudio de los morfismos desde el punto

\oldpage[I]{8}
de vista de la dimensión (cap.~IV), que se transcribe sin cambios notables al marco de los esquemas.
Conviene señalar, por otra parte, que la noción de «esquemas de anillos locales», introducida por Chevalley, se presta naturalmente a una extensión de la geometría algebraica (sin tener, no obstante, toda la flexibilidad y la generalidad que pretendemos darle aquí); sobre las relaciones entre esta noción y nuestra teoría, véase el capítulo primero, \S~8.
M.~Nagata ha desarrollado tal extensión en una serie de trabajos~\cite{I-9} que contienen numerosos resultados particulares sobre la geometría algebraica sobre anillos de Dedekind\footnote{Entre los trabajos próximos a nuestro punto de vista en geometría algebraica, señalemos el importante trabajo de E.~Kähler~\cite{I-22} y una nota reciente de Chow e Igusa~\cite{I-3}, que retoman, en el marco de la teoría de Nagata-Chevalley, ciertos resultados de (FAC) y dan también una fórmula de Künneth.}.

\bigskip
\begin{center}*\quad*\quad*\end{center}
\bigskip
